%%%%%%%%%%%%%%%%%%%%%%%%%%%%%%%%%%%%%%%%%%%%%%%%%%%%%%%%%%%%%%%%%%%%%%%
% agents-prefer-agents — TAIGR @ ICML 2026 submission (anonymous)
%
% Compile (requires texlive):
%   cd paper && pdflatex paper && bibtex paper && pdflatex paper && pdflatex paper
%
% Target length: 2 main pages + unlimited appendix.
%%%%%%%%%%%%%%%%%%%%%%%%%%%%%%%%%%%%%%%%%%%%%%%%%%%%%%%%%%%%%%%%%%%%%%%

\documentclass{article}

\usepackage{microtype}
\usepackage{graphicx}
\usepackage{booktabs}
\usepackage{hyperref}
\usepackage{amsmath}
\usepackage{amssymb}
\usepackage[capitalize,noabbrev]{cleveref}

% Initial blind version.
\usepackage{icml2026}

\icmltitlerunning{Precursors to Gradual Disempowerment on GitHub}

\begin{document}

\twocolumn[
\icmltitle{Precursors to gradual disempowerment: \\
  AI-to-AI interaction chains are growing in GitHub pull requests}

\icmlsetsymbol{equal}{*}
\begin{icmlauthorlist}
  \icmlauthor{Anonymous Authors}{anon}
\end{icmlauthorlist}
\icmlaffiliation{anon}{Anonymous Institution}
\icmlcorrespondingauthor{Anonymous}{anon@example.com}
\icmlkeywords{AI agents, gradual disempowerment, GitHub, code review, technical AI governance}

\vskip 0.3in
]

\printAffiliationsAndNotice{}

%%%%%%%%%%%%%%%%%%%%%%%%%%%%%%%%%%%%%%%%%%%%%%%%%%%%%%%%%%%%%%%%%%%%%%
\begin{abstract}
The gradual-disempowerment scenario of \citet{kulveit2025gradual} predicts that delegated markets can decouple from human preferences as AI agents replace humans in their operation. The concept is theoretical; its \emph{precursors} -- observable, measurable patterns that would precede such a drift -- have not been systematically quantified in the wild. We offer a \textbf{first illustration} of how to do so, on a concrete ecosystem: GitHub pull-request review. We define a taxonomy of nine precursors, build an event-level detector for AI-agent activity, and illustrate the measurement of two precursors on 84,725{} pull requests from 1,000{} popular GitHub repositories (April~2025--March~2026). In this sample, AI-agent participation grew 3.5{}$\times$ in one year (from 6.5\%{} to 22.5\%{} of all PRs; contiguous AI$\to$AI event chains of length~$\geq$5 grew from 0.3\%{} to 4.8\%{}). A first within-PR test -- of the 161{} PRs where an AI bot issued an \emph{approving} review, humans co-approved 5/24 (20.8\%){} of 24{} AI-authored PRs versus 46/137 (33.6\%){} of 137{} human-authored ($z=-1.24$, two-sided $p=0.22${}) -- finds no statistically significant AI-AI bias, not yet adjusted for code quality. We release the pipeline and treat these measurements as an initial probe; the contribution is the \emph{measurability} of precursors more than any particular point estimate.
\end{abstract}

\section{Introduction}
\label{sec:intro}

``Gradual disempowerment''~\citep{kulveit2025gradual} names a macro failure in which economic, cultural, and political systems drift from human preferences as AI increasingly replaces human labour within them. The concept is theoretical and its individual mechanisms are hard to study in isolation: they require longitudinal observation of real multi-agent systems at scale. Agent deployment is now observable at that scale~\citep{appel2025economic,stein2026mcp,shao2025workbank}, sharpening the risk of overreliance~\citep{ibrahim2025overreliance}. Prior work on \emph{AI-AI bias} has shown that LLMs prefer LLM-generated content over human content in single-turn choice experiments---GPT-4 favoured LLM-written product ads 89\% of the time vs.\ a 36\% human baseline, and LLM abstracts 78\% vs.\ 61\%~\citep{laurito2025aibias}; LLM evaluators recognise and favour their own generations~\citep{panickssery2024evaluators}. Whether analogous biases manifest \emph{in the wild}, in multi-agent systems that actually decide real-world outcomes, is an open empirical question, and so is the prior question of whether any of the gradual-disempowerment \emph{precursors} are measurable at all outside the lab.

This paper is a \textbf{first illustration} of precursor measurement in a concrete ecosystem: GitHub pull-request review. We choose GitHub because real AI coding agents act in it (Claude~Code, GitHub~Copilot coding agent, Devin, Cursor, Aider, Cline), distinct AI code-review agents (CodeRabbit, Cubic, Gemini~Code~Assist, Copilot~PR~Reviewer, Greptile) leave machine-readable traces~\citep{ghaleb2026fingerprinting,staufer2026aiagentindex}, merge outcomes are discrete, and a full year (April~2025--March~2026) covers the period in which coding-agent adoption became mainstream. Concurrent work reports high agent-PR merge rates~\citep{watanabe2025agentic,li2026aidev,li2025teammates,logicstar2025wild} alongside persistent quality debt~\citep{agarwal2026ides,ehsani2026fail,huang2026morecode}; none of it frames the measurement in disempowerment-precursor terms.

We offer three illustrative contributions. (i)~A \emph{first taxonomy} of nine precursors distinguished from the situational (within-conversation) failure patterns studied in prior work~\citep{kulveit2025gradual,sharma2026whos}. (ii)~A \emph{first open measurement pipeline} that classifies GitHub PR \emph{events} (rather than whole repositories as in \citet{ghaleb2026fingerprinting}) as human-, AI-assisted-, or AI-bot-attributed, usable as a standing monitor. (iii)~A \emph{first illustrative measurement} of two precursors on 84,725{} PRs from 1,000{} repositories, including -- to our knowledge -- the first within-PR comparison of AI-bot and human approvals on the same code; prior empirical work compares disjoint pools of PRs~\citep{zhong2026synergy,chowdhury2026claims}. Given small numerators (the key cell has only 24{} PRs), we emphasise the \emph{feasibility} of measurement; our effect-size estimates are illustrative and associational, and intended as a starting point for larger-scale tracking, not as conclusive findings.

\section{A taxonomy of precursors}
\label{sec:taxonomy}

\begin{table}[h]
\centering\footnotesize
\caption{Nine precursors to gradual disempowerment in multi-agent ecosystems. ``Measured'' = this paper; others are future work.}
\label{tab:taxonomy}
\begin{tabular}{@{}lll@{}}
\toprule
Precursor & Short definition & Here? \\
\midrule
Chain length & length of uninterrupted AI$\leftrightarrow$AI events & \textbf{yes} \\
AI review preference & AI approves AI output vs.\ human & \textbf{yes} \\
Model-family self-pref. & same family preferred over other AI & future \\
AI-entity preference & AI defers to AI over routable human & future \\
Output homogeneity & AI outputs cluster tighter than humans & future \\
Anti-human collusion & coordinated exclusion of a human actor & future \\
Self-empowerment & AI expands its own tokens/permissions & future \\
Overconfidence & AI proceeds where human would ask & future \\
Erratic under stress & VendingBench-style failure modes & future \\
\bottomrule
\end{tabular}
\end{table}

\emph{Precursors} are observable, measurable patterns; they are not themselves disempowerment. \citet{sharma2026whos} focus on situational (within-conversation) failures; our taxonomy is ecosystem-level and is designed to be tracked from public logs over time. The table is a starting draft rather than a settled classification -- we invite extensions -- and it complements within-system MAS failure taxonomies~\citep{cemri2025mas} and autonomy-as-user-role frameworks~\citep{cihon2025autonomy,feng2025autonomylevels}.

\section{Data and methods}
\label{sec:methods}

\paragraph{Repository universe.} Public GitHub repositories plausibly gaining $\geq$1{,}000 stars in 2025 and actively pushed within our analysis window, capped at the top 1,000{} by a recency-weighted activity score; not forks, not archives. Non-goal: GitHub-wide representativeness.

\paragraph{PR universe.} All PRs \emph{created} between 2025-04-01 and 2026-03-31 in those repositories (84,725{} PRs), fetched via GitHub GraphQL with \texttt{orderBy:{CREATED\_AT,DESC}}. We record metadata, commits, reviews, review-thread comments, issue comments, and merge/close timeline events.

\paragraph{Event-level AI classification.} Adapting \citet{ghaleb2026fingerprinting}'s heuristics to event granularity, we classify each actor--event pair as \textbf{AI-bot} (account on our 20-entry allowlist: \texttt{copilot-swe-agent[bot]}, \texttt{devin-ai-integration[bot]}, \texttt{cursor-agent}, \texttt{claude[bot]}, \texttt{coderabbitai}, \texttt{cubic-dev-ai}, \texttt{gemini-code-assist}, \texttt{copilot-pull-request-reviewer}, etc.\ -- see Appendix~C), \textbf{AI-assisted} (human account whose payload contains a \texttt{Co-Authored-By} trailer to an AI tool or a direct handle mention), or \textbf{human}. Confidence is \emph{high} when a bot account or co-author trailer matches, \emph{low} otherwise. Maintenance bots (Dependabot, Renovate, CI) are excluded. Primary analyses use high-confidence AI only.

\paragraph{Chain length (RQ1).} Within a PR, order attributed events by timestamp. An \emph{AI$\to$AI chain} is a maximal contiguous subsequence of events with \texttt{actor\_type=AI-bot} at high confidence. The per-PR metric is the longest such chain; we report the monthly share of PRs with chain $\geq$1, $\geq$2, $\geq$5.

\paragraph{Self-preference diff-in-diff (RQ2).} For each PR we assign \texttt{author\_type}~$\in\{$AI,human$\}$ (AI if the PR's author login is on our allowlist OR any commit bears an AI co-author trailer) and \texttt{reviewer\_type}~$\in\{$AI,human$\}$ (AI if an approving review came from our allowlist). Tie-breaking is conservative: AI wins ties, biasing \emph{against} finding AI-AI preference. Restricting to PRs with at least one approving review, the cell merge rate (within 30 days of opening) is $\mathrm{mr}[a,r]$; the diff-in-diff is
\[
    \mathrm{DiD} = (\mathrm{mr}_{AA} - \mathrm{mr}_{Ah}) - (\mathrm{mr}_{hA} - \mathrm{mr}_{hh}).
\]
Confidence intervals use a cluster bootstrap by repository (1{,}000 draws).

\paragraph{Scope: authors and reviewers, not mergers.} We measure AI participation through authoring and reviewing, not merging. No merges in our universe are attributed to AI-bot accounts on our allowlist, so ``AI merger'' is not an informative dimension in this window; we discuss the limitations of that framing---including the fact that AI can still act \emph{through} a human merger---in the Discussion.

\section{Results}
\label{sec:results}

\begin{figure*}[t]
    \centering
    \includegraphics[width=0.98\linewidth]{figures/figure1_ai_participation.pdf}
    \caption{Headline: monthly share of PRs with any AI-agent event, and with contiguous AI$\to$AI chains of length $\geq$2 and $\geq$5, across 1,000{} popular GitHub repositories (Apr 2025--Mar 2026). Vertical dotted lines mark approximate public-availability dates of several coding and code-review agents (reader context; no causal claim).}
    \label{fig:headline}
\end{figure*}

\paragraph{Chain length (RQ1): participation is growing.} Figure~\ref{fig:headline} illustrates monthly agent-participation rates at three intensity levels in this sample. AI-agent participation rose from 6.5\%{} of PRs in April~2025 to 22.5\%{} in March~2026 (3.5{}$\times$). The share of PRs containing an \emph{uninterrupted} chain of at least two AI-to-AI events grew from 2.7\%{} to 16.3\%{}; chains of length $\geq$5 from 0.3\%{} to 4.8\%{}. The p95 longest-chain length doubled from 1 to 4 events. The trend is monotonic across twelve months. We read this as illustrative evidence that the \emph{chain-length} precursor is measurable and moving; a single high-activity slice of GitHub cannot establish that AI-agent activity is accelerating ecosystem-wide.

\paragraph{Self-preference (RQ2): a first look.} Table~\ref{tab:did} reports the $2\times2$ merge rates on PRs with approving reviews. The diff-in-diff is +0.2 percentage points (95\% CI -31.2 to +25.1){}; the 95\% cluster-bootstrap CI straddles zero. The table is best read as illustrative: AI code-review tools are not obviously preferentially approving AI-authored PRs, and AI-reviewed PRs in general have lower merge rates than human-reviewed ones. Only two AI bots in our data (Cubic, CodeRabbit) routinely issue approving reviews, which constrains how much weight this first contrast can carry.

\paragraph{Within-PR comparison (Fig.~\ref{fig:withinpr}).} A tighter test holds the PR (and implicitly, code quality) fixed. AI code-review tools in our sample overwhelmingly issue COMMENTED reviews rather than APPROVED ones: of 17,169{} AI-bot review events, only 172{} ($\approx$1.0\%{}) are approvals, and these come almost entirely from Cubic~AI (n=85) and CodeRabbit (n=19); Copilot~PR~Reviewer, Gemini~Code~Assist, and Greptile post critiques but almost never approve. The rejection channel is equally thin (AI bots issued \texttt{changes\_requested} on only 11 PRs). Using ``AI reviewer'' strictly to mean ``AI bot that issued an APPROVED review'', we ask whether a human reviewer and an AI approver, looking at the same PR, agree more readily when the PR is AI-authored. Of the 161{} PRs where an AI bot approved, 24{} (14.9\%{}) were AI-authored; a human co-approved 5/24 (20.8\%){} of those, versus 46/137 (33.6\%){} of the 137{} human-authored AI-approved PRs. The point-estimate gap (+12.7{}~pp) is directionally consistent with a small AI-AI bias, but $z=-1.24$, two-sided $p=0.22${} and the 95\% CIs overlap substantially. \textbf{With a key cell of only 24{} PRs, we treat this as an illustration of the test, not a finding}: we demonstrate that a within-PR AI-AI-bias measurement is feasible on public GitHub data, and report the first such contrast. A larger approving-agent population is required before the direction can be read as a substantive claim.

\begin{figure}[t]
    \centering
    \includegraphics[width=\linewidth]{figures/figure2_withinpr.pdf}
    \caption{Human co-approval rate among the 161{} PRs where an AI bot issued an APPROVED review. Red: AI-authored (n=24{}). Blue: human-authored (n=137{}). Point estimates differ by +12.7{} percentage points (directionally consistent with AI reviewers approving AI code humans disagree with more often), but the 95\% CIs overlap ($z=-1.24$, two-sided $p=0.22${}) --- no statistically significant replication of the \citet{laurito2025aibias} single-turn AI-AI bias in this narrow channel. Error bars: Wilson 95\% CIs. Not adjusted for code quality.}
    \label{fig:withinpr}
\end{figure}

\begin{table}[t]
\centering\footnotesize
\caption{Merge rates by author$\times$reviewer type (approved PRs only). Cluster-bootstrap DiD = +0.2{} percentage points (95\% CI -31.2{} to +25.1{}).}
\label{tab:did}
\begin{tabular}{@{}lrr@{}}
\toprule
 & AI reviewer & human reviewer \\
\midrule
AI author & 75.0\%{} (n=24{}) & 87.8\%{} (n=1,325{}) \\
human author & 75.4\%{} (n=134{}) & 88.4\%{} (n=25,746{}) \\
\bottomrule
\end{tabular}
\end{table}

\section{Discussion and call to action}
\label{sec:discuss}

\paragraph{What the illustration shows.} As a demonstration of measurability, three things hold in this sample. Chain length is easy to track and is moving in the direction predicted by disempowerment concerns. The cross-cell DiD is estimable but structurally under-powered because so few AI bots issue approving reviews. The within-PR contrast is the tightest test currently feasible on public GitHub data; with 24{} AI-authored, AI-approved PRs it is under-powered for an effect size at the single-digit-pp scale. A larger population of approving AI agents --- or a longer window --- would convert this illustration into a test with meaningful power. AI-mediated peer review~\citep{thakkar2025iclr} and social sycophancy~\citep{cheng2025elephant} suggest that AI can shape human judgments without overt preference, consistent with overreliance risk~\citep{ibrahim2025anthropomorphic}; our measurement apparatus would let such effects, if they emerge, be tracked over time rather than discovered retrospectively.

\paragraph{Call to action.} Treat precursor tracking as a first-order empirical task rather than a by-product of capability evaluation. If this pipeline can be built from public logs in a few hundred lines of Python, similar measurements should become standing dashboards: platforms (GitHub, Stack Overflow, npm, package registries) should publish monthly AI-agent activity reports; model providers should include cross-family self-preference in post-training evaluations; sectoral regulators should require the same for regulated agent markets (finance, healthcare). The raw signal is already in public logs; this paper is a first illustration that it can be extracted and tracked.

\paragraph{Limitations.} We flag seven main limitations; full discussion in Appendix~B.
(i)~\emph{Selection.} Our 500 repositories are the top of a high-activity slice of GitHub (median stars 53{,}583; see Appendix~E for the exact pipeline). We do not claim representativeness of GitHub-wide PR behaviour.
(ii)~\emph{Detection false-negatives.} Event-level heuristic detection cannot recover deliberately unattributed agent activity. Commercial coding agents reportedly include modes that suppress attribution in outward artefacts; the headline growth figure is therefore a lower bound.
(iii)~\emph{Detection false-positives.} Our 20-entry bot allowlist (Appendix~C) may include accounts that are AI-\emph{assisted} but not AI-\emph{agent} in the sense defined in §\ref{sec:intro}. We accept this trade-off because false positives move PRs from `human' into `AI', \emph{reducing} the measured author-$\times$-reviewer contrast and biasing \emph{against} our null finding.
(iv)~\emph{Thin AI$\times$AI cell for RQ2.} Only two bots in our data (Cubic AI, CodeRabbit) routinely post \emph{approving} reviews; the AI-author $\times$ AI-reviewer cell holds only 24{} approved PRs even across 500 repos and twelve months. The DiD CI is therefore wide by construction, not by chance.
(v)~\emph{Confounding by tool announcements.} Several coding- and review-agent GAs happened within our window (vertical lines in Fig.~\ref{fig:headline}). We make no causal attribution; the chain-length trend is associational.
(vi)~\emph{Cluster-bootstrap assumptions.} Resampling repositories assumes exchangeability conditional on the $2\times 2$ cell; a few very active repositories contribute disproportionately.
(vii)~\emph{Merger attribution.} We focus on AI \emph{reviews}, not AI \emph{merges}. No merges in our universe are attributed to AI-bot accounts on our allowlist, so direct ``AI merger'' activity is currently zero. A human still clicks the merge button, but AI may still act \emph{through} that human---e.g., when an AI review's ``approved'' state is what gates the maintainer's decision. Our design cannot distinguish a maintainer who merged \emph{because of} an AI approval from one who merged in spite of it; a fuller causal account of merge gating is future work.
We make no causal claims in either direction.

%%%%%%%%%%%%%%%%%%%%%%%%%%%%%%%%%%%%%%%%%%%%%%%%%%%%%%%%%%%%%%%%%%%%%%
\bibliography{references}
\bibliographystyle{icml2026}

%%%%%%%%%%%%%%%%%%%%%%%%%%%%%%%%%%%%%%%%%%%%%%%%%%%%%%%%%%%%%%%%%%%%%%
\newpage\appendix
\onecolumn
\section{Example pull requests}
\label{app:examples}
\input{appendix/examples.tex}

\section{Robustness and threats to validity}
\label{app:threats}
% Appendix B — robustness and threats to validity (expanded from §Limitations)

This appendix expands the limitations listed in the main paper.

\paragraph{B.1 Selection into the repository universe.}
We restrict to \PLACEHOLDERPRACTIVE{} critical OS software repositories ranked by OpenSSF \texttt{criticality\_score} (\Cref{app:repo-selection}). The resulting sample has median stars \PLACEHOLDERSTARMED{} (range \PLACEHOLDERSTARMIN{}--\PLACEHOLDERSTARMAX{}) and a criticality-score range of \PLACEHOLDERSCOREMIN{}--\PLACEHOLDERSCOREMAX{} --- it skews toward load-bearing, multi-org, heavily-cross-referenced open-source infrastructure (language toolchains, runtimes, foundation projects) rather than uniform popularity. Effects we estimate apply to ``critical OS software repositories''. We do not generalise to GitHub as a whole. In particular, chain-length growth in lower-criticality repositories may be smaller or larger than what we report, and we do not have data to say which.

\paragraph{B.2 Detection false-negatives: ``undercover'' agents.}
Our event-level detection (\S\ref{sec:methods}, \Cref{app:allowlist}) relies on two signals: bot-account logins, and \texttt{Co-Authored-By} trailers naming an AI tool. Commercial coding agents have been reported to include configuration options that suppress attribution in outward artefacts~---~e.g., an ``undercover mode'' reportedly present in at least one major coding agent's system prompt, revealed via a source-code leak in early 2026. PRs produced by such configurations appear as ordinary human activity to our detector. Consequently the measured AI-agent participation share in Figure~\ref{fig:headline} is a \emph{lower bound}. This does not overturn the paper's main claim (participation is growing) but tightens its interpretation: the true level is above what we measure, by an unknown amount.

\paragraph{B.3 Thin AI-approval slice in the within-PR test.}
Among our \PLACEHOLDERNREPOS{} repositories and \PLACEHOLDERNPRS{} PRs, only \PLACEHOLDERWITHINN{} carry an approving review from an AI bot, and only \PLACEHOLDERWITHINAIN{} of those are also AI-authored. This is not a sampling accident: fewer than ten bots in our allowlist ever emit an ``APPROVED'' review state, and only two (Cubic, CodeRabbit) do so routinely. The two-proportion $z$-test on the 24 vs.\ 137 cells has limited power. We report it as a null on explicit self-preference, with the understanding that larger-sample follow-ups are needed to distinguish a small positive effect from zero at the $\pm 5$~pp level.

\paragraph{B.4 Merger attribution: authors and reviewers, not mergers.}
We initially considered ``AI merger'' as a third dimension (author $\times$ reviewer $\times$ merger). No merges in our universe are attributed to AI bot accounts on our allowlist, so the merger dimension collapses to ``human'' on every observed PR. This is a substantive limitation: a human clicks the merge button on every merged PR in our window, but AI can still act \emph{through} that human---most obviously when an AI review's ``approved'' state is what gates a maintainer's decision to merge. Our observational design cannot separate a maintainer who merged \emph{because of} an AI approval from one who merged in spite of it. We also caution that our detector excludes merges performed under a maintainer's personal-access token by an AI app or merge-queue automation acting on their behalf. Such merges would appear as human in both the REST \texttt{merged\_by} field and the GraphQL timeline actor, and therefore cannot be recovered from the public event log alone.

\paragraph{B.5 Event timestamp ordering and chain definition.}
AI$\to$AI chains are defined on events sorted by timestamp within a PR. GitHub timestamps are accurate to seconds. Ties are broken by the API's default order (commit $<$ review $<$ review-comment $<$ issue-comment). Alternative tie-breaking does not change the head of the distribution meaningfully (mean and $p_{95}$ shift by $<$0.1 events). The tail (max chain = 13) is insensitive.

\paragraph{B.6 Reproducibility.}
Every number in the paper is derived from one of \texttt{data/pr\_events.parquet} (one row per actor/event) or \texttt{data/chains.parquet} (one row per PR, chain aggregates). These are produced by \texttt{scripts/03\_classify\_prs.py} and \texttt{05\_compute\_chains.py} respectively. The within-PR two-proportion test is produced by \texttt{scripts/06c\_within\_pr.py}. The upstream PR-event data (\texttt{data/prs/*.jsonl}) is collected by \texttt{scripts/02\_fetch\_prs.py} against the \PLACEHOLDERNREPOS{} repositories in \texttt{data/repos.json}, which itself is produced by \texttt{scripts/01\_build\_repo\_list.py} (\Cref{app:repo-selection} details this step). Given a valid GitHub token, the full pipeline reproduces the paper's numbers up to GitHub's own indexing non-determinism.


\section{Bot and AI-tool allowlists}
\label{app:allowlist}
\input{appendix/allowlist.tex}

\section{Full regression}
\label{app:reg}
% Appendix D — full logit regression with controls

We estimate
\[
    \mathrm{logit}(\text{merged in 30d}) = \alpha + \beta_1 \mathbf{1}\{\text{AI author}\} + \beta_2 \mathbf{1}\{\text{AI reviewer}\} + \beta_3 \mathbf{1}\{\text{AI author}\}\cdot\mathbf{1}\{\text{AI reviewer}\} + \gamma \cdot \mathbf{x} + \eta_\text{repo} + \eta_\text{week} + \varepsilon,
\]
where $\mathbf{x}$ contains $\log(1+\text{additions}+\text{deletions})$, author tenure (days since first commit to the repo), and a CI-status indicator where available. Standard errors are clustered by repository.

\PLACEHOLDERREGRESSIONTABLE

The coefficient of interest is $\beta_3$, the interaction of AI-author and AI-reviewer. \PLACEHOLDERREGRESSIONINTERPRETATION{}


\section{Repository-selection pipeline (replication)}
\label{app:repo-selection}
% Appendix --- Sample-selection pipeline

\paragraph{Overview.} The repository universe is selected via the OpenSSF
\emph{criticality\_score} \citep{ossf-criticality}.
The score is a published, replicable
ranking of GitHub projects by their structural importance to the open-source
ecosystem; using it directly addresses the well-known critique that
star-based selection over-weights short-lived viral projects relative to
load-bearing infrastructure
\citep{borges2018github, kalliamvakou2014promises, munaiah2017curating}.
We pin a published criticality snapshot, then filter and enrich it via the
GitHub REST API.

\paragraph{The OpenSSF criticality score.} For each project, the score is a
weighted arithmetic mean over normalised signals
\citep{ossf-criticality}. The signals used by the default \texttt{all.csv}
configuration are:
\begin{itemize}
\item \texttt{created\_since}: months since first commit;
\item \texttt{updated\_since}: months since last commit (negatively weighted);
\item \texttt{contributor\_count}: distinct authors over the project's history;
\item \texttt{org\_count}: distinct organisations among contributors;
\item \texttt{commit\_frequency}: average commits per week (last year);
\item \texttt{recent\_release\_count}: GitHub releases in the last year;
\item \texttt{updated\_issues\_count} and \texttt{closed\_issues\_count}: issues touched in the last 90 days;
\item \texttt{issue\_comment\_frequency}: comments per issue;
\item \texttt{github\_mention\_count}: cross-repo references (in commit messages and other repos' code).
\end{itemize}
The score lies in $(0, 1]$; higher is more critical. The OpenSSF
publishes a global score CSV to a public Google Cloud Storage bucket
(\texttt{gs://ossf-criticality-score/}) with a new snapshot every two
weeks. We pin a single snapshot (Step~1 below) so that re-running the
pipeline reproduces the same universe.

\paragraph{Step 1 --- Pin the criticality snapshot.} We use the
\PLACEHOLDERSNAPSHOT{} snapshot (\texttt{all.csv} variant, no deps.dev
dependent counts), which contains \PLACEHOLDERSNAPSHOTROWS{} scored
projects.

\paragraph{Step 2 --- Filter \& sort by score.} From the
\PLACEHOLDERSNAPSHOTROWS{} rows we keep those with a numeric
\texttt{default\_score} and a \texttt{repo.url} starting with
\texttt{https://github.com/}. After filtering: \PLACEHOLDERSCOREDROWS{}
rows. We sort descending by \texttt{default\_score}, breaking ties by
\texttt{repo.url} alphabetically for stability across reruns.

\paragraph{Step 3 --- Candidate cap (over-sample).} We take the top
\PLACEHOLDERCANDCAP{} candidates by score. The over-sample (vs.\ the
final \PLACEHOLDERFINALCAP{} target) absorbs drop-out in the next two
filtering steps: repos that no longer exist on GitHub, were renamed to a
private slug, or that are forks/archived/inactive.

\paragraph{Step 4 --- GitHub enrichment.} For each candidate we issue
\texttt{GET /repos/\{owner\}/\{name\}} (REST API, async, 8-way bounded
concurrency). The result populates the standard repo-metadata fields
(stars, forks, \texttt{pushed\_at}, license, default branch, etc.). HTTP
404 (deleted, renamed, or now-private) drops the row. After enrichment we
reject:
\begin{itemize}
\item \texttt{fork == true} (forks contribute no independent PR activity);
\item \texttt{archived == true};
\item \texttt{disabled == true};
\item \texttt{pushed\_at} strictly before \texttt{2025-04-01} (no
in-window PRs);
\item zero in-window PRs on GitHub (canonical development happens off GitHub --- LKML, git.apache.org, git.eclipse.org, etc.).
\end{itemize}

\paragraph{Step 5 --- Final cap.} Of the remaining \PLACEHOLDERPOSTACTIVITY{}, we keep the top
\PLACEHOLDERFINALCAP{} by criticality score as the final repo list.
Each row carries the GitHub-API fields plus the score, the within-snapshot
rank, the snapshot date, and the raw OpenSSF signal columns (so any
downstream re-weighting is possible without re-fetching the CSV).

\paragraph{Sample characteristics.}
\begin{itemize}
\item \PLACEHOLDERFINALCAP{} repositories selected from
\PLACEHOLDERSNAPSHOTROWS{} scored candidates $\to$ \PLACEHOLDERCANDCAP{}
top-by-score $\to$ \PLACEHOLDERENRICHOK{} successful GitHub enrichments
$\to$ \PLACEHOLDERPOSTACTIVITY{} after activity filter $\to$
\PLACEHOLDERFINALCAP{} after cap.
\item Criticality-score range: max \PLACEHOLDERSCOREMAX{}, min \PLACEHOLDERSCOREMIN{}.
\item Stars: min \PLACEHOLDERSTARMIN{}, median \PLACEHOLDERSTARMED{}, max \PLACEHOLDERSTARMAX{}.
\item Top languages: \PLACEHOLDERLANGTOP{}.
\end{itemize}

\paragraph{Analysed PRs.}
\begin{itemize}
\item Total PRs analysed: 3,109,979 (mean 311.0, median 93 per repo; p90 843, max 7,321).
\item Per-repo cap of \PLACEHOLDERPRSCAP{} PRs: 2 of \PLACEHOLDERFINALCAP{} repos (0.0\%) saturate the cap; for these the sampled rows are weighted toward the centre of the window by week-uniform subsampling rather than the most-recent months. All other repos have their full in-window PR history analysed.
\end{itemize}

\paragraph{Known selection biases.} The criticality score is itself an
opinionated metric: it favours older projects with multi-org contributor
bases, recent commit activity, frequent releases, and inbound cross-repo
references. Three biases inherited from this:
\begin{itemize}
\item Library-style projects score higher than application-style projects
of comparable importance because they accumulate inbound mentions from
their dependents' commit messages.
\item Pure-binary or assets-only repositories with no commits-as-history
(e.g.\ Wikipedia mirrors) are under-scored.
\item The score is computed over GitHub-visible signals only; non-GitHub
mirrors of important projects (e.g.\ many Apache and Eclipse Foundation
projects whose canonical home is elsewhere) are absent or under-counted.
\end{itemize}


\section{Supplementary figures}
\label{app:extra-figures}
\input{appendix/extra_figures.tex}

\end{document}
